O gerenciamento lógico da infraestrutura de rede compreende o planejamento, implementação e manutenção da configuração da camada 2 (switches e VLANs), da camada 3 (roteamento) e da camada 4 (firewall). Esta atividade é estratégica para garantir isolamento adequado entre serviços e desempenho da rede. As atividades constantes relacionadas a esta responsabilidade são: 

\begin{itemize} 
    \item Planejamento e criação de VLANs para segmentação lógica da rede; 
    \item Configuração de switches;
    \item Monitoramento do estado físico dos switches (temperatura, ventilação e ruído anormal) com SNMP;
    \item Verificação periódica de logs de hardware e alertas críticos;
    \item Acompanhamento do consumo de CPU, memória e utilização de portas;
    \item Inspeção de fontes e módulos de uplink;
    \item Testes de portas de uplink com falhas ou intermitência;
    \item Identificação e substituição de cabos de uplink defeituosos;
    \item Atualização e controle de inventário (patrimonial);
    \item Acionamento de garantia junto ao fabricante em caso de falha de hardware;
    \item Apoio na resposta a incidentes de segurança relacionados à rede; 
    \item Planejamento de endereçamento IP e organização de sub-redes;
    \item Definição de políticas de comunicação entre VLANs (no HP são as Zonas);
    \item Criação e manutenção de regras de firewall para controle de acesso entre redes internas e externas;
    \item Revisão periódica das regras de firewall e políticas de acesso;
    \item Implementação de NAT quando necessário;
    \item Monitoramento de tráfego para identificação de gargalos ou comportamentos anômalos;
\end{itemize}
